\documentclass[10pt,a4paper]{amsart}
\usepackage[margin=19mm]{geometry}
\usepackage{amsmath,amssymb,mathtools}
\usepackage[scr=rsfs,cal=euler]{mathalfa}
\usepackage{xfrac}
\usepackage[unicode,hidelinks,bookmarks=false]{hyperref}
\newcommand{\ie}{i.e.\ }
\newcommand{\cf}{cf.\ }
\newcommand{\resp}{respectively\ }
\newcommand{\Subsection}{\subsection}
\providecommand{\nobreakdash}{\nobreak\hbox{-}}
\setlength{\emergencystretch}{2em}
\multlinegap0pt
\usepackage{bm}
\usepackage{bbm}
\usepackage{dsfont}
\usepackage{xspace}
\usepackage{cancel}
\usepackage{mathtools}
\usepackage{amsthm}
\makeatletter
\@ifundefined{thm}{\newtheorem{thm}{Thm}}{}
\@ifundefined{prop}{\newtheorem{prop}[thm]{Proposition}}{}
\makeatother
\title[Integral Kirwan Surjectivity]{Integral Kirwan Surjectivity}
\author{Daniel Pomerleano, Constantin Teleman}
\date{Source: arXiv:2410.06197v2 (CC BY 4.0)}
\begin{document}
\maketitle

\noindent\emph{Source and licence.} Integral Kirwan Surjectivity. Version arXiv:2410.06197v2 (CC BY 4.0), distributed under the Creative Commons Attribution 4.0 International licence, as recorded in the arXiv licence metadata for this exact version and shown on the abstract page https://arxiv.org/abs/2410.06197v2. Full rights notice: licenses/arxiv-2410.06197-CC-BY-4.0.txt.\par\smallskip

\noindent\textit{Editorial scope.} Counted unit: the Prop whose source label is \texttt{prop:prMoravacircle} in the article (source0.tex lines 691--693, source label \texttt{prop:prMoravacircle}), delivered together with the article's own complete original proof of it (source0.tex lines 710--761, one \texttt{proof} environment of 52 lines). No mathematics is written, repaired or re-derived: statement and proof are the article's own text line for line. The proof is detached from the statement in the source: the article states the result at lines 691--693 and prints its proof at lines 710--761, and the proof environment's own header names the statement, so the association is the article's own. Required context retained uncounted: eq:restrictioneuler (lines 512--514); eq:degenerationm1 (lines 686--688); eq:Morseinequalities (lines 698--700). Every definition, display and label of those lines is retained unchanged. Omitted from this package and not redistributed: 1--511, 515--685, 689--690, 694--697, 701--709, 762--1298. Nothing inside a retained excerpt is omitted or altered: every retained body is delivered exactly as the article's lines are written, so no omission receipt of any kind accompanies this package. Citations: the retained excerpts carry no citation command at all, so no bibliography entry of the article is transcribed and no reference is redistributed. Reference adaptation: none was needed and none was made; every reference inside the delivered unit resolves inside it, so no unresolved reference or citation is produced. Printed slips kept literal: no printed word, symbol, hypothesis or number of the counted statement or of its proof was altered. Layout and class: the article's own class is \texttt{amsart} with packages including amsmath, amssymb, amsfonts, amstext, verbatim, amsthm, mathtools, setspace, enumerate, prettyref, fullpage, tikz-cd, stmaryrd; the wrapper is the lane's pinned \texttt{amsart} editorial wrapper at 10pt on A4, which restates the theorem-like environments the retained text needs and adds the article's own macro definitions that the retained text needs (none), with extra wrapper packages (bm, bbm, dsfont, xspace, cancel, mathtools). The delivered numbering is the unit's own. Assets: no retained span contains a figure environment, a graphics-inclusion command, a PDF-page inclusion, a table or a TikZ picture; no image file is copied into this package, the package declares no asset dependency and no image file is redistributed with it. Licence: version arXiv:2410.06197v2 of this article is distributed under the Creative Commons Attribution 4.0 International licence, as recorded in the arXiv licence metadata for this exact version and shown on the abstract page https://arxiv.org/abs/2410.06197v2. A full rights notice is delivered at licenses/arxiv-2410.06197-CC-BY-4.0.txt. Characters: every retained excerpt is pure ASCII, so the delivered engine, XeLaTeX with its default Latin Modern text font, reports no missing character.
\begin{align} \label{eq:restrictioneuler} 
e_{S^1\times H}(V)= u^kx + m \in E^*(N_{S^1 \times H}),  
\end{align}  

\begin{align} 
\label{eq:degenerationm1} E_2^{s,t}= H^s(BT, K_{p^r}(n)^t(M)) \Rightarrow K_{p^r}(n)^{s+t}(M_T), 
\end{align}

\begin{align} \label{eq:Morseinequalities} 
|K_{p^r}(n)^*(F)| \geq  |K_{p^r}(n)^*(M)| 
\end{align} 

\begin{prop} \label{prop:prMoravacircle} 
The  spectral sequence \eqref{eq:degenerationm1} collapses at $E_2$.    
\end{prop}

\begin{proof}[Proof of Proposition~\ref{prop:prMoravacircle}]
Start with the special case $T=S^1$ and let $F_{S^1,m}$ and $M_{S^1,m}$ denote the restrictions to 
$\mathbb{CP}^m$ of the respective Borel fibrations to $BS^1=\mathbb{CP}^\infty$. Then,  
\begin{align} 
\label{eq:Leraycardinalitybound} 
|K_{p^r}(n)^*(M_{S^1,m})| \leq |K_{p^r}(n)^*(M)|^{m+1}. 
\end{align} 
Write $K_{p^r}(n)^*(F_{S^1,m})=\bigoplus_{j=0}^m K_{p^r}(n)^*(F)u^j$,  where 
$u \in K_{p^r}(n)^2(\mathbb{C}P^m)$ is the generator.  From the form \eqref{eq:restrictioneuler} of the 
Euler class, we extract an integer $k \geq 1$, depending on the $p$-exponents in the weights of 
the action of $S^1$ on $\nu$, but \emph{independent of~$m$}, such that 
\begin{align} 
\label{eq:kernelbound} 
\operatorname{ker}(\iota_m^* \iota_{m*}) \bigcap 
	\bigoplus_{j=0}^{m-k} K_{p^r}(n)^*(F)u^j = 0
\end{align}  
for the restricted $\iota_m$ over $\mathbb{CP}^m$; this is because $xu^k$ is a leading term of the Euler 
class in the $\mathfrak{m}$-filtration. The projection of 
$\operatorname{ker}(\iota_m^* \iota_{m*})$ 
onto the complementary space spanned by powers $u^j$ with 
$m-k< j\le m$ is injective. It follows that
\begin{align}
\label{eq:Eulerinjectivitybound} |K_{p^r}(n)^*(F)|^{m+1}-|K_{p^r}(n)^*(F)|^{k+1} \leq |K_{p^r}(n)^*(M_{S^1,m})|.  
\end{align} 
Combining \eqref{eq:Leraycardinalitybound} with~\eqref{eq:Eulerinjectivitybound} for large $m$ 
shows that we must have equality in \eqref{eq:Morseinequalities}. 

Suppose now that we find a non-trivial differential on some page of the 
sequence~\eqref{eq:degenerationm1}.  Multiplying it  by  powers of the hyperplane class in
$H^2(\mathbb{CP}^m)$ would show that the difference  
\[
|K_{p^r}(n)^*(M)|^{m+1}-|K_{p^r}(n)^*(M_{S^1,m})|
\] 
grows unboundedly with $m$,  contradicting the $k$-bound in \eqref{eq:Eulerinjectivitybound}. 
This concludes the proof for $T=S^1$.  

For a general $T$, we proceed by induction on the rank, splitting  $T=S^1 \times H$. 
Choose  finite-dimensional approximations 
$BT_m =  \mathbb{CP}^m\times BH_m\cong \mathbb{CP}^m \times \cdots \times \mathbb{CP}^m$, 
with corresponding $M_{T,m}$. We claim  the $E_2$ degeneration of all Leray sequences  
\begin{align} \label{eq:moravaprdegm} 
E_2^{s,t}= H^s(BT_m, K_{p^r}(n)^t(M)) \Rightarrow K_{p^r}(n)^{s+t}(M_{T,m}).
\end{align} 
Restricting from $BH$ to $BH_m$, our induction hypothesis implies the collapse of \eqref{eq:moravaprdegm} with $H$ replacing $T$.  Since $M_{H,m}$ is a Hamiltonian $S^1$-space, 
we conclude that \eqref{eq:moravaprdegm} also degenerates with $T$.  
From this,  we deduce the surjectivity of the restriction-to-fiber map for all 
$m$, 
\begin{align} \label{eq:degenerationmorava} 
K_{p^r}(n)^*(M_{T,m}) \to K_{p^r}(n)^*(M)
\end{align} 
ensuring the absence of differentials and collapse in~\eqref{eq:degenerationm1}. 
\end{proof} 

\end{document}
